\documentclass[12pt]{article}
\title{Bayesian Statistics}
\date{\today}
\usepackage{amsmath}

\setlength{\parindent}{0pt} % don't indent the line
\begin{document}
\maketitle

\section{Posterior Distribution}
$\theta$ is the parameter. $D$ is data. $P(\theta)$ is prior distribution. $P(\theta|D)$ is posterior distribution.\\

Probability of parameter and data is defined as: $P(\theta \cap D) = P(\theta)P(D|\theta) = P(D)P(\theta|D)$ \\

Posterior distribution can be calculated by rearrangement: $P(\theta|D) = \frac{P(\theta)P(D|\theta)}{P(D)}$ \\

\section{Markov chain Monte Carlo}
MCMC estimates the posterior, $P(\theta|D)$, without estimating $P(D)$. \\

The goal of MCMC is to find a range with the parameter. \\

MCMC process:
\begin{enumerate}
    \item Choose a starting point in parameter space based on the prior .
    \item Move from one point to another based on: $P(\frac{A \rightarrow B}{B \rightarrow A} \propto \frac{Post(B)}{Post(A)})$, where A and B are two points in parameter space.
    \item The process will spend more time around a certain point.
    \item The process stops after the posterior distribution is not changing (converges).
    \item The program will output a list of all points explored during the process.
    \item A range with the true parameter is be picked from the list.
\end{enumerate}



\end{document}